\begin{figure}
    \begin{center}
    \includegraphics[width=1.0\linewidth]{Figures/CongressionalBills_Estimation/LLM_On_LHS/fig03 Bias Density LLM.jpeg}
    \caption{Bias Density of $\beta^\text{LLM}$. Proxy on the LHS.}
    % \label{fig:enter-label}
    \end{center}
    Notes: 
    The bias estimates were calculated as the difference between the LLM-based coefficient $\beta^\text{LLM}$ and the human-based coefficient $\beta^\star$, using a regression of the form $Y_t = \alpha + \beta V$, where $t \in$ \{Health; Banking, Finance, and Domestic Commerce; Defense; Government Operations; Public Lands and Water Management; Other\} represents major topics and $V \in \{\text{Democrat}; \text{Senate}; \text{DW1}\}$ are the covariates. 
    For $\beta^\star$, $Y_t^\text{Human} := 1\{Y^\text{Human} = t\}$, while for $\beta^\text{LLM}$, $Y_{t,m,p}^\text{LLM} := 1\{Y_{m,p}^\text{LLM} = t\}$, with $m \in \{\text{GPT-3.5}, \text{GPT-4o}\}$ denoting the LLM model and $p \in \{1, \ldots, 12\}$ representing the prompt. 
    The human-based coefficient $\beta^\star$ was computed using the full 10,000 bills, whereas the LLM-based coefficient $\beta^\text{LLM}$ was calculated as the average of $N = 1,000$ simulation runs, with each run’s estimate based on a sample of 5,000 bills. 
    The bias estimates are independent of the validation proportion $\eta$, though they were generated in the experiment run with $\eta = 5\%$.
\end{figure}